% !TeX root = ../../Book5.tex
% 纲要标题：### 6.2 Specht 模
% 纲要位置：计划纲要.md，第 3260 行。

\section{Specht 模}
\label{b5:sec:0602}

% TODO: 按以下纲要要点撰写本节。

% 纲要内容（仅作写作计划，尚非教材正文）：
%
% * 表类置换模、tabloid（按行置换的表类）与 polytabloid（列反对称化的表类线性组合）；原“多重表”不再作为 polytabloid 的译名，multitableau 另作术语辨析
% * Young 对称化子的特征零构造
% * Specht 模的基本性质
% * 明确 Young 表、按行置换取等价类的 tabloid 与列反对称化所得 polytabloid 的区别，统一表模和 Specht 模的符号
% * 标准 polytabloid 基与 Garnir 关系的证明集中到第6.3节；Young 对称化子构造与该模型的关系同时交代
%

待撰写。
